\begin{table}[!htbp]\centering\footnotesize
\caption{Real data: Type A worlds versus covariate-preserving worlds}\label{tab:v_covA}
\begin{tabular}{lrr}
\toprule
 & Type A (within sector) & Covariate-preserving \\
\midrule
Worlds & 1,000 & 1,000 \\
Largest null H1 $t$: median / 95th pct.\ / max & 7.4 / 8.6 / 9.8 & 7.2 / 8.0 / 9.3 \\
Search-adjusted $p$, H1 (real $t=12.1$) & 0.001 & 0.001 \\
Search-adjusted $p$, full package & 0.056 & 0.137 \\
Potemkin rate, U (\%) & 26.5 & 40.3 \\
Potemkin rate, U-free (\%) & 93.4 & 100.0 \\
\bottomrule
\end{tabular}
\begin{minipage}{0.97\linewidth}\vspace{2pt}\scriptsize\textit{Notes.} Covariate-preserving worlds keep each exposure proxy's within-sector projection on twelve predetermined traits (the ten moderators, pre-period size and hiring growth; within-sector $R^2$ 53\% for the main proxy) and permute only the residual within sector, so trait-driven outcome trends are present in every null world.\end{minipage}
\end{table}